%! Function Fundamentals: Notation, Domain, and Range — Unit 2, Lesson 2.1 — Exit Ticket
\documentclass[10pt]{article}
\usepackage{atda-article}
\usepackage{atda-boxes}
\newcommand{\TallMath}[1]{$\displaystyle #1\rule[-1.4em]{0pt}{3.2em}$}
\setlength{\workrowsep}{8pt}

\begin{document}

\pageheader{Unit 2, Lesson 2.1}{Exit Ticket}
% NAMESTRIP: no \namedateperiod here. The student writes their name once, on the
% cover this component is stapled behind. See references/conventions.md.

\vspace{0.15in}

\begin{notesbox}{Both Directions --- On Your Own}
\small
Notes closed. Three items. A swimming pool is drained at a steady rate. $W$ is the water left in
the pool, in gallons, $t$ hours after the pump was switched on.

\begin{center}
\begin{tikzpicture}
\begin{axis}[
    width=9.8cm, height=4.4cm,
    xlabel={$t$ (hours after the pump starts)}, ylabel={$W$ (gallons)},
    xmin=0, xmax=8.6, ymin=0, ymax=2650,
    xtick={0,1,2,3,4,5,6,7,8}, ytick={0,600,1200,1800,2400},
    grid=both, grid style={linegray!45}, axis lines=left,
    tick label style={font=\scriptsize}, label style={font=\scriptsize}]
  \addplot[royal, very thick, domain=0:8, samples=2]{2400-300*x};
\end{axis}
\end{tikzpicture}
\end{center}

\begin{enumerate}[leftmargin=1.6em, itemsep=8pt, topsep=4pt]

\item Read $W(2)$ off the graph. Write it as a point, then as a sentence about the pool.

\writelines{2}

\item Now the other direction: find the value of $t$ for which $W(t)=600$, and say what that
number means for whoever is draining the pool.

\writelines{2}

\item State the domain and the range in this situation, with units, in inequality notation
\emph{and} in interval notation. Then explain why $t=10$ is not in the domain, even though
$2400-300(10)$ is a number you can compute.

\writelines{3}

\end{enumerate}
\end{notesbox}

\end{document}
